\section{2019级工科数学分析下 B卷补考}

\subsection{资料来源}
\begin{itemize}
  \item 试题：2019 级工科数学分析（二）B 卷补考用，已导出页面图校正公式。
  \item 未发现配套答案或解析文本；下列答案与解析为据题面补写。
\end{itemize}

\subsection{试题}

\paragraph*{试卷信息}
诚信应考，考试作弊将带来严重后果！

华南理工大学本科生期末考试 2019--2020--2 学期《工科数学分析（二）》B卷。

注意事项：开考前请将密封线内各项信息填写清楚；所有答案请直接答在试卷上；考试形式：闭卷；本试卷共（五）大题，满分100分，考试时间120分钟。评阅教师请在试卷袋上评阅栏签名。

\paragraph*{一、填空题（共5小题，每题3分，共15分）}
\begin{enumerate}
  \item 函数 \(z=x^2+2y^2\) 在点 \((1,1)\) 沿该点梯度方向的方向导数为 \underline{\hspace{5cm}}。
  \item 曲线 \(x=t,\ y=t^2,\ z=t^3\) 在点 \((1,1,1)\) 的法平面方程为 \underline{\hspace{5cm}}。
  \item 函数 \(z=xy\) 在条件 \(x^2+y^2=1\) 下的极大值为 \underline{\hspace{5cm}}。
  \item 级数 \(\displaystyle \sum_{n=1}^{\infty}\frac{x^n}{n\cdot4^n}\) 的和函数为 \underline{\hspace{5cm}}。
  \item 设 \(S(x)\) 为 \(f(x)=\mathrm{e}^{\cos x}\)，\(x\in[0,\pi]\)，展成的以 \(2\pi\) 为周期的余弦级数的和函数，则 \(S(0)=\) \underline{\hspace{5cm}}。
\end{enumerate}

\paragraph*{二、计算题（共5小题，每题9分，共45分）}
\begin{enumerate}
  \item 求微分方程
  \[
    \frac{\mathrm{d}y}{\mathrm{d}x}=\frac{2x-y+2}{x+2y-4}
  \]
  的通解。

  \item 假设隐函数存在定理的条件都满足，且 \(y(t)\) 由方程组
  \[
    \begin{cases}
      y=f(x,t),\\
      x=g(y,t)
    \end{cases}
  \]
  确定，求 \(\dfrac{\mathrm{d}y}{\mathrm{d}t}\)。

  \item 计算
  \[
    \iiint_\Omega (x^2+y^2+z^2)\,\mathrm{d}x\mathrm{d}y\mathrm{d}z,
  \]
  其中 \(\Omega\) 是曲面 \(z=x^2+y^2\) 与 \(z=1-x^2-y^2\) 所围成的区域。

  \item 计算曲面积分
  \[
    \iint_\Sigma z\,\mathrm{d}S,
  \]
  其中 \(\Sigma\) 是球面 \(x^2+y^2+z^2=4\) 被平面 \(z=\sqrt2\) 截出的顶部。

  \item 计算
  \[
    I=\iint_{\Sigma^+}\frac{\cos(\vec r,\vec n)}{r^2}\,\mathrm{d}S,
  \]
  其中 \(\Sigma^+\) 为不经过 \(P(a,b,c)\) 的光滑闭曲面，方向向外，\(\Sigma\) 不包含点 \(P(a,b,c)\)。
  \(\vec n\) 为曲面 \(\Sigma\) 上任一点 \(M(x,y,z)\) 处的外法向量的单位向量，\(\cos(\vec r,\vec n)\) 表示两向量 \(\vec r,\vec n\) 夹角的余弦，且
  \[
    \vec r=\overrightarrow{MP}=(a-x,b-y,c-z),\qquad
    r=|\vec r|=\sqrt{(a-x)^2+(b-y)^2+(c-z)^2}.
  \]
\end{enumerate}

\paragraph*{三、解答下列各题（共2小题，每题8分，共16分）}
\begin{enumerate}
  \item 试确定 \(\alpha\) 的值，使得函数
  \[
    f(x,y)=
    \begin{cases}
      \dfrac{x^\alpha y^\alpha}{x^2+y^2}, & (x,y)\ne(0,0),\\
      0, & (x,y)=(0,0)
    \end{cases}
  \]
  在点 \((0,0)\) 处可微。
  \item 将 \(\ln(1-3x+2x^2)\) 展成麦克劳林级数。
\end{enumerate}

\paragraph*{四、证明题（共2小题，每题8分，共16分）}
\begin{enumerate}
  \item 设 \(f\) 为连续函数，\(C\) 为分段光滑的简单闭曲线，证明
  \[
    \oint_C f(x^2-y^2)(x\,\mathrm{d}x-y\,\mathrm{d}y)=0.
  \]
  \item 证明函数项级数
  \[
    \sum_{n=1}^{\infty}\frac{(-1)^n(1-\mathrm{e}^{-nx})}{n^2+x^2}
  \]
  在 \([0,+\infty)\) 上一致收敛。
\end{enumerate}

\paragraph*{五、应用题（共1题，每题8分，共8分）}
求点 \(\left(2,2,\dfrac12\right)\) 到抛物面 \(z=x^2+y^2\) 的最小距离。

\subsection{答案与解析}
\paragraph*{一、填空题}
\begin{enumerate}
  \item \(\nabla z(1,1)=(2,4)\)，沿梯度方向的方向导数为
  \[
    \|\nabla z(1,1)\|=2\sqrt5.
  \]
  \item 曲线切向量为 \((1,2t,3t^2)\)，在 \(t=1\) 处为 \((1,2,3)\)。法平面为
  \[
    (x-1)+2(y-1)+3(z-1)=0,
  \]
  即 \(x+2y+3z-6=0\)。
  \item 由 \(2xy\le x^2+y^2=1\)，得极大值为
  \[
    \frac12.
  \]
  \item
  \[
    \sum_{n=1}^{\infty}\frac{x^n}{n4^n}
    =\sum_{n=1}^{\infty}\frac{(x/4)^n}{n}
    =-\ln\left(1-\frac x4\right),\qquad -4\le x<4.
  \]
  \item 余弦级数在端点取偶延拓的函数值，故
  \[
    S(0)=f(0)=\mathrm{e}.
  \]
\end{enumerate}

\paragraph*{二、计算题}
\begin{enumerate}
  \item 作平移
  \[
    X=x,\qquad Y=y-2.
  \]
  方程化为
  \[
    \frac{\mathrm{d}Y}{\mathrm{d}X}=\frac{2X-Y}{X+2Y}.
  \]
  写成微分形式：
  \[
    (2X-Y)\,\mathrm{d}X-(X+2Y)\,\mathrm{d}Y=0.
  \]
  这是恰当方程，积分得
  \[
    X^2-XY-Y^2=C.
  \]
  因而原方程通解为
  \[
    x^2-x(y-2)-(y-2)^2=C.
  \]

  \item 对
  \[
    y=f(x,t),\qquad x=g(y,t)
  \]
  关于 \(t\) 求导，得
  \[
    y'=f_xx'+f_t,\qquad x'=g_yy'+g_t.
  \]
  消去 \(x'\)，得到
  \[
    \frac{\mathrm{d}y}{\mathrm{d}t}
    =\frac{f_t+f_xg_t}{1-f_xg_y}.
  \]

  \item 用柱坐标。区域为
  \[
    0\le r\le \frac1{\sqrt2},\qquad r^2\le z\le 1-r^2,\qquad 0\le\theta\le2\pi.
  \]
  所以
  \[
    I=2\pi\int_0^{1/\sqrt2}\int_{r^2}^{1-r^2}(r^2+z^2)r\,\mathrm{d}z\,\mathrm{d}r
    =\frac{11\pi}{96}.
  \]

  \item 在球面 \(x^2+y^2+z^2=4\) 上取球坐标 \(z=2\cos\varphi\)，顶部对应
  \(0\le\varphi\le\pi/4\)。于是
  \[
    \iint_\Sigma z\,\mathrm{d}S
    =\int_0^{2\pi}\int_0^{\pi/4}2\cos\varphi\cdot4\sin\varphi\,\mathrm{d}\varphi\,\mathrm{d}\theta
    =4\pi.
  \]

  \item 被积函数可看作向量场
  \[
    \mathbf F(M)=\frac{\overrightarrow{MP}}{|\overrightarrow{MP}|^3}
  \]
  通过 \(\Sigma\) 的通量。若点 \(P\) 在 \(\Sigma\) 所围区域外，则 \(\mathbf F\) 在该区域内无奇点且散度为 \(0\)，由 Gauss 公式可得
  \[
    I=0.
  \]
  若点 \(P\) 在 \(\Sigma\) 所围区域内，由于 \(\vec r=\overrightarrow{MP}=P-M\) 与标准径向场方向相反，通量为
  \[
    I=-4\pi.
  \]
\end{enumerate}

\paragraph*{三、解答题}
\begin{enumerate}
  \item 由极坐标估计，
  \[
    |f(x,y)|\le C r^{2\alpha-2}.
  \]
  若在原点可微且导数为零，需要
  \[
    \frac{|f(x,y)|}{r}\le C r^{2\alpha-3}\to0.
  \]
  因此
  \[
    \alpha>\frac32.
  \]
  此时 \(f_x(0,0)=f_y(0,0)=0\)，且 \(f(x,y)=o(\sqrt{x^2+y^2})\)，故函数在原点可微。

  \item 因为
  \[
    1-3x+2x^2=(1-x)(1-2x),
  \]
  所以
  \[
    \ln(1-3x+2x^2)=\ln(1-x)+\ln(1-2x)
    =-\sum_{n=1}^{\infty}\frac{1+2^n}{n}x^n,\qquad |x|<\frac12.
  \]
\end{enumerate}

\paragraph*{四、证明题}
\begin{enumerate}
  \item 取原函数 \(F'(u)=f(u)\)。由于
  \[
    \mathrm{d}(x^2-y^2)=2x\,\mathrm{d}x-2y\,\mathrm{d}y,
  \]
  故
  \[
    f(x^2-y^2)(x\,\mathrm{d}x-y\,\mathrm{d}y)
    =\frac12\,\mathrm{d}F(x^2-y^2).
  \]
  闭曲线上的全微分积分为 \(0\)，因此结论成立。

  \item 对 \(x\in[0,+\infty)\)，有
  \[
    \left|\frac{(-1)^n(1-\mathrm{e}^{-nx})}{n^2+x^2}\right|
    \le\frac1{n^2}.
  \]
  因为 \(\sum_{n=1}^{\infty}1/n^2\) 收敛，由 Weierstrass 判别法，原函数项级数在 \([0,+\infty)\) 上一致收敛。
\end{enumerate}

\paragraph*{五、应用题}
设抛物面上点为 \((x,y,x^2+y^2)\)。距离平方为
\[
  D=(x-2)^2+(y-2)^2+\left(x^2+y^2-\frac12\right)^2.
\]
由对称性，最近点满足 \(x=y=s\)。于是
\[
  D(s)=2(s-2)^2+\left(2s^2-\frac12\right)^2.
\]
令
\[
  D'(s)=8(2s^3-1)=0,
\]
得 \(s=2^{-1/3}\)。最小距离平方为
\[
  D_{\min}^2=\frac{33}{4}-3\sqrt[3]{4}.
\]
因此最小距离为
\[
  \sqrt{\frac{33}{4}-3\sqrt[3]{4}}.
\]
